\documentclass{article}


% if you need to pass options to natbib, use, e.g.:
%     \PassOptionsToPackage{numbers, compress}{natbib}
% before loading neurips_2023


% ready for submission
\usepackage{neurips_2023}


% to compile a preprint version, e.g., for submission to arXiv, add add the
% [preprint] option:
%     \usepackage[preprint]{neurips_2023}


% to compile a camera-ready version, add the [final] option, e.g.:
%     \usepackage[final]{neurips_2023}


% to avoid loading the natbib package, add option nonatbib:
%    \usepackage[nonatbib]{neurips_2023}


\usepackage[utf8]{inputenc} % allow utf-8 input
\usepackage[T1]{fontenc}    % use 8-bit T1 fonts
\usepackage{hyperref}       % hyperlinks
\usepackage{url}            % simple URL typesetting
\usepackage{booktabs}       % professional-quality tables
\usepackage{amsfonts}       % blackboard math symbols
\usepackage{nicefrac}       % compact symbols for 1/2, etc.
\usepackage{microtype}      % microtypography
\usepackage{xcolor}         % colors


\title{Formatting Instructions For NeurIPS 2023}


% The \author macro works with any number of authors. There are two commands
% used to separate the names and addresses of multiple authors: \And and \AND.
%
% Using \And between authors leaves it to LaTeX to determine where to break the
% lines. Using \AND forces a line break at that point. So, if LaTeX puts 3 of 4
% authors names on the first line, and the last on the second line, try using
% \AND instead of \And before the third author name.


\author{%
  David S.~Hippocampus\thanks{Use footnote for providing further information
    about author (webpage, alternative address)---\emph{not} for acknowledging
    funding agencies.} \\
  Department of Computer Science\\
  Cranberry-Lemon University\\
  Pittsburgh, PA 15213 \\
  \texttt{hippo@cs.cranberry-lemon.edu} \\
  % examples of more authors
  % \And
  % Coauthor \\
  % Affiliation \\
  % Address \\
  % \texttt{email} \\
  % \AND
  % Coauthor \\
  % Affiliation \\
  % Address \\
  % \texttt{email} \\
  % \And
  % Coauthor \\
  % Affiliation \\
  % Address \\
  % \texttt{email} \\
  % \And
  % Coauthor \\
  % Affiliation \\
  % Address \\
  % \texttt{email} \\
}


\begin{document}


\maketitle


\begin{abstract}
  The abstract paragraph should be indented \nicefrac{1}{2}~inch (3~picas) on
  both the left- and right-hand margins. Use 10~point type, with a vertical
  spacing (leading) of 11~points.  The word \textbf{Abstract} must be centered,
  bold, and in point size 12. Two line spaces precede the abstract. The abstract
  must be limited to one paragraph.
\end{abstract}

\section{Introduction}

This report examines the application of the Expectimax algorithm to the game 2048 as a representative example
of decision-making under uncertainty. Unlike adversarial environments, 2048 combines deterministic player actions
with stochastic outcomes through random tile generation. This structure makes Expectimax more suitable than Minimax,
as it evaluates decisions based on expected outcomes rather than worst-case scenarios. Additionally, the game
requires strategic planning, which can be effectively approximated using heuristic evaluation within a depth-limited
search. Therefore, 2048 provides a concise case for analyzing Expectimax in a probabilistic setting.

\subsection*{What is the Expectimax algorithm?}

The Expectimax algorithm is a decision-making method designed for environments involving both deterministic choices
and probabilistic events. It extends Minimax by introducing chance nodes, which represent random outcomes and compute
expected values based on their probabilities.
\\
\\
In 2048, player moves are modeled as Max nodes, while random tile generation (2 with 90\% probability and 4 with 10\%) is represented as chance nodes. The algorithm explores future states up to a fixed depth, selecting moves that maximize expected outcomes. Due to computational constraints, a heuristic evaluation function (e.g., snake pattern, empty tiles, smoothness) is used to estimate non-terminal states. This combination enables efficient decision-making in a stochastic environment.


\section{Method}


Please read the instructions below carefully and follow them faithfully. \textbf{Important:} This year the checklist will be submitted separately from the main paper in OpenReview, please review it well ahead of the submission deadline: \url{https://neurips.cc/public/guides/PaperChecklist}.


\subsection{Implementation}

\subsubsection{Board State and Transitions}

Each state of 2048 game is represented as a $4 \times 4$ board, where each tile contains a number $2^k$ with
non-negative $k$. A value of 0 denotes an empty tile. At each step, the agent chooses one of four actions:
left, right, up, or down. After each move, a new tile is added to a randomly selected empty cell.
The tile value is 2 with probability 0.9 and 4 with probability 0.1. When two tiles of value $x$ are merged,
they produce a tile of value $2x$, and the score is increased by $2x$. The total game score is accumulated
as the sum of all merge operations. A state is considered terminal when no valid moves remain.

\subsubsection{Expectimax Search}

To handle the stochastic nature of the game, we employ the Expectimax Search algorithm to select actions.
The game is modeled as a search tree consisting of two types of nodes: max nodes and chance nodes.
\\
\\
Max nodes correspond to the agent's decision points. From a given state, the agent considers all valid moves
and selects the action that maximizes the expected value of future states. Chance nodes represent the generation
of new tiles. From a chance node, all possible tile placements are considered, each weighted by its probability.
The heuristic value of a chance node is a weighted mean of heuristics of all its children nodes.
\\
\\
Formally, the heuristic of a node is calculated recursively. At a max node, the heuristic is the maximum over
its children where $s'$ is the successor state after applying action $a$.

$$h(s) = \max_{a \in A(s)} h(s')$$

At a chance node, the heuristic value is computed as the expected value over all possible outcomes.

$$h(s) = \sum_{s'} P(s'|s) \cdot h(s')$$

Due to the exponential growth of the search tree, we impose a depth limit on the Expectimax search.
When the maximum depth is reached or no further moves are available, the board state is evaluated directly
using the heuristic function $h(s)$.

\subsection{Heuristic Function Design}

In order to effectively evaluate non-terminal states in the Expectimax search algorithm for the 2048 games,
we designed heuristic functions that approximate the desirability of a current state. 
\\
\\
We represent a board state as a $4 \times 4$ grid $S$, where each tile value is denoted as $s_{i,j}$.
The overall heuristic is composed of three components: snake structure, empty cells, and smoothness, each capturing
complementary properties of a desirable board.

\subsubsection{Snake Heuristic}

$$H_{snake}(S) = \sum_{i=1}^{4} \sum_{j=1}^{4} w_{i,j} \cdot s_{i,j}$$

The snake heuristic evaluates whether high-value tiles are arranged in a structured, monotonic order.
This is implemented using a predefined weight matrix $W$, which assigns larger weights to positions along a snake
path starting from a corner of the board. As a result, it enforces a global ordering and facilitates future merge
operations.

\subsubsection{Empty Cells Heuristic}

$$H_{empty}(S) = 1000 \cdot \sum_{i, j} 1(s_{i,j} = 0)$$

The empty cells heuristic measures the number of unoccupied positions on the board. A higher number of empty cells
implies more possible future moves and a lower risk of reaching a terminal state. The scaling factor is included to
make this term more comparable in magnitude to the other heuristic values used in evaluation.

\subsubsection{Smoothness Heuristic}

$$H_{smooth}(S) = -\sum_{(i,j)} \sum_{(k,l)\in N(i,j)} |s_{i,j} - s_{k,l}|$$

The smoothness heuristic measures the closeness of magnitude between adjacent tiles on the grid.
It penalizes large differences in value between neighboring tiles, assigning higher scores to grids where
neighboring tiles have similar magnitudes. In the implementation, the smoothness term is computed only when 
$s_{i, j} \neq 0$, and differences are measured only with the right and lower neighboring tiles.

\subsubsection{Combined Heuristic}

$$H(S) = H_{snake}(S) + H_{empty}(S) + H_{smooth}(S)$$

The combined heuristic is defined as the sum of the three components. An additive formulation is used to ensure
that all three criteria are considered simultaneously. Unlike multiplication or max-based approaches, summation
allows each component to contribute proportionally and enables meaningful trade-offs between them when evaluating
a board state.

\section{Experiment}

\subsection{Variables and Results}

\begin{itemize}
  \item Controlled Variables
  \begin{itemize}
    \item Game environment ($4\times4$ board size, game rules, and termination condition)
    \item Tile generation mechanism (probability distribution: 2 with 90\%, 4 with 10\%)
    \item Decision-making framework (Expectimax algorithm structure, including Max and Chance nodes)
    \item Action space (four directional moves: up, down, left, right)
    \item Number of trials per configuration (25 runs)
  \end{itemize}
  \item Independent Variables
  \begin{itemize}
    \item Heuristic function (snake, empty, smooth, combined)
    \item Search depth (1, 2, 3)
  \end{itemize}
  \item Dependent Variables
  \begin{itemize}
    \item Average score (Avg)
    \item Maximum score (Max)
    \item Minimum score (Min)
    \item Standard deviation (Std)
    \item Maximum tile achieved (MaxTile)
    \item Success rates ($\ge1024$, $\ge2048$)
  \end{itemize}
\end{itemize}

\begin{table}[h]
\centering
\begin{tabular}{llrrrrrrr}
\toprule
Heuristic & Depth & Avg & Max & Min & Std & Max Tile & 2048 Rate & 1024 Rate \\
\midrule

snake & 1 & 3299 & 10440 & 756 & 1948 & 1024 & 0/25 & 1/25 \\
snake & 2 & 3208 & 7140 & 740 & 1562 & 512 & 0/25 & 0/25 \\
snake & 3 & 35599 & 111280 & 3156 & 22195 & 8192 & 20/25 & 22/25 \\

\midrule
empty & 1 & 2990 & 6724 & 788 & 1356 & 512 & 0/25 & 0/25 \\
empty & 2 & 2990 & 6724 & 788 & 1356 & 512 & 0/25 & 0/25 \\
empty & 3 & 9075 & 15568 & 3352 & 3792 & 1024 & 0/25 & 10/25 \\

\midrule
smooth & 1 & 1974 & 4724 & 552 & 987 & 512 & 0/25 & 0/25 \\
smooth & 2 & 2579 & 7012 & 660 & 1502 & 512 & 0/25 & 0/25 \\
smooth & 3 & 6633 & 15980 & 1740 & 3909 & 1024 & 0/25 & 5/25 \\

\midrule
combined & 1 & 4570 & 9084 & 1992 & 2030 & 512 & 0/25 & 0/25 \\
combined & 2 & 6296 & 14368 & 1848 & 3170 & 1024 & 0/25 & 3/25 \\
combined & 3 & 18220 & 58688 & 5448 & 12095 & 4096 & 7/25 & 19/25 \\

\bottomrule
\end{tabular}
\caption{Performance of Expectimax with Different Heuristics and Depths (25 Trials each)}
\end{table}

\subsection{Comparison with Human Performance}

To evaluate the performance of our model, we compare its gameplay results with those of human players.
Human scores were collected from 12 gameplay runs. We used combined heuristic with depth 3 for comparison.

\begin{table}[h]
\centering
\begin{tabular}{lrrrr}
\toprule
Model & Avg Score & Max Score & Max Tile & $\geq$2048 \\
\midrule
Human & 6303 & 16944 & 1024 & 0\% \\
Agent (snake depth 3) & 35599 & 111280 & 8192 & 80\% \\
\bottomrule
\end{tabular}
\caption{Comparison of human performance and agent performance}
\end{table}

The model significantly outperforms human players in both score and maximum tile achieved.
While human players did not reach 2048, the model achieved 2048 20 out of 25 times, even with limited depth of 3.

\section{Discussion}

\subsection*{Analysis}

The results from Figure 1 show that search depth has the greatest impact on Expectimax performance.
Increasing the depth from 1 to 3 significantly improved average scores across all heuristics
(e.g., snake: $3,299 \rightarrow 35,599$), highlighting the importance of deeper lookahead.
\\
\\
Heuristic choice had limited effect at shallow depths but became critical at depth 3.
The snake heuristic achieved the best performance, with the highest average score (35,599) and the highest success 
rate in reaching 2048 (20/25).
\\
\\
The snake heuristic proved to be the dominant contributor to agent performance. When tuning the combined heuristic,
reducing the weight of the empty cell and smoothness components while increasing the weight of the snake component
consistently improved results, suggesting that the positional strategy encoded by the snake pattern is the most
valuable signal for 2048 play. This is further supported by the fact that snake alone outperforms all other individual
heuristics and the combined heuristic on average score at every depth tested. However, the combined heuristic
demonstrated greater robustness at depth 3, achieving a lower standard deviation (12,095 vs 22,195) and a higher
minimum score (5,448 vs 3,156) compared to snake alone. This suggests that while the empty cell and smoothness
components do not improve peak performance, they act as stabilisers, reducing the frequency of poor outcomes and
producing more consistent play across games. The snake heuristic alone is therefore sufficient for maximising average
performance, but the combined heuristic may be preferable in settings where consistency is valued over peak score.

\subsection*{Improvement}

Reinforcement learning agents learn strategy through millions of self-play games, allowing them to discover
patterns and long-term planning that a hand-crafted heuristic like ours cannot capture. Rather than relying on
a human-designed evaluation function, they develop an intuition for what makes a board good purely from experience.
This makes them far less dependent on design choices like heuristic weighting or search depth, which as our results
show are significant bottlenecks for expectimax. 
\end{document}
